% GENERATED FILE. Edit canonical lessons, then run npm run generate:books.
% canonical-source-hash: fnv1a64:6292a43b81a3ba5c
% canonical-lessons: FA-C538-mostajer, FA-C538-maykroviv, FA-C538-lavazem-e-khanegi, FA-C538-garm-kardan, FA-C538-sard-shodan

\chapter{Renting and Appliances}
\label{ch:fa-renting-and-appliances}
\hlchaptermodality{538}

\begin{chapteropening}
\textbf{By the end of this chapter:} \emph{I can talk about renting and household appliances.}

Complete the last lesson of chapter 538: I can talk about renting and household appliances.
\end{chapteropening}

\section[\texorpdfstring{\fa{مستاجر} (mosta'jer)}{مستاجر (mosta'jer)}]{\fa{مستاجر} (mosta'jer) — a tenant}

\label{lesson:FA-C538-mostajer}



\begin{quote}
Before the new one: say the Persian for to break down, then the Persian for a landlord.
\end{quote}

\subsection*{You'll want to know: \fa{مستاجر}}
\textbf{\fa{مستاجر}} (\emph{mosta'jer}) — \textquotedblleft{}a tenant\textquotedblright{}.

\begin{quote}
\fa{مستاجر} \fa{طبقه} \fa{بالا} \fa{خیلی} \fa{ساکت} \fa{است}.
\end{quote}

\begin{quote}
\emph{mosta'jer-e tabaqe-ye bâlâ kheyli sâkat ast.}
\end{quote}

\begin{quote}
The tenant upstairs is very quiet.
\end{quote}

Plural \emph{mosta'jerân}. The opposite is \emph{sâhebkhâne}.

\begin{cousinweb}
\emph{mosta'jer} is a loan, usually traced to Arabic \emph{mustaʾjir}, 'one who rents', from the root of \emph{ejâre}, 'rent'. EtymDB has no entry for it.
\end{cousinweb}

\begin{grammarlens}[title={What you've built}]
One more word for telling, explaining and describing at B1.
\end{grammarlens}

\subsection*{Your turn}
\begin{itemize}
\raggedright
  \item \emph{Say it:} \emph{mosta'jer}
  \item \emph{Say it:} \emph{mosta'jer}, once more
  \item \emph{Say it:} a sentence of your own with \emph{mosta'jer}
  \item \emph{Recall:} say the Persian for to break down, then the Persian for a landlord, then say \emph{mosta'jer} again
\end{itemize}

\begin{tcolorbox}[breakable,colback=teal!4,colframe=teal!35!black,title={Before you move on}]
What does \fa{مستاجر} mean? (\textquotedblleft{}A tenant\textquotedblright{}.) Say it once more.
\end{tcolorbox}
\section[\texorpdfstring{\fa{مایکروویو} (mâykroviv)}{مایکروویو (mâykroviv)}]{\fa{مایکروویو} (mâykroviv) — a microwave (oven)}

\label{lesson:FA-C538-maykroviv}



\begin{quote}
Before the new one: say the Persian for a landlord, then the Persian for a tenant.
\end{quote}

\subsection*{You'll want to know: \fa{مایکروویو}}
\textbf{\fa{مایکروویو}} (\emph{mâykroviv}) — \textquotedblleft{}a microwave (oven)\textquotedblright{}.

\begin{quote}
\fa{مایکروویو} \fa{ما} \fa{خراب} \fa{شد}.
\end{quote}

\begin{quote}
\emph{mâykroviv-e mâ kharâb shod.}
\end{quote}

\begin{quote}
Our microwave broke down.
\end{quote}

Plural \emph{mâykrovivhâ}.

\begin{cousinweb}
\emph{mâykroviv} is a modern loan, usually traced to English \emph{microwave}. EtymDB has no entry for it.
\end{cousinweb}

\begin{grammarlens}[title={What you've built}]
One more word for telling, explaining and describing at B1.
\end{grammarlens}

\subsection*{Your turn}
\begin{itemize}
\raggedright
  \item \emph{Say it:} \emph{mâykroviv}
  \item \emph{Say it:} \emph{mâykroviv}, once more
  \item \emph{Say it:} a sentence of your own with \emph{mâykroviv}
  \item \emph{Recall:} say the Persian for a landlord, then the Persian for a tenant, then say \emph{mâykroviv} again
\end{itemize}

\begin{tcolorbox}[breakable,colback=teal!4,colframe=teal!35!black,title={Before you move on}]
What does \fa{مایکروویو} mean? (\textquotedblleft{}A microwave (oven)\textquotedblright{}.) Say it once more.
\end{tcolorbox}
\section[\texorpdfstring{\fa{لوازم} \fa{خانگی} (lavâzem-e khânegi)}{لوازم خانگی (lavâzem-e khânegi)}]{\fa{لوازم} \fa{خانگی} (lavâzem-e khânegi) — household appliances}

\label{lesson:FA-C538-lavazem-e-khanegi}



\begin{quote}
Before the new one: say the Persian for a tenant, then the Persian for a microwave.
\end{quote}

\subsection*{You'll want to know: \fa{لوازم} \fa{خانگی}}
\textbf{\fa{لوازم} \fa{خانگی}} (\emph{lavâzem-e khânegi}) — \textquotedblleft{}household appliances\textquotedblright{}.

\begin{quote}
\fa{این} \fa{فروشگاه} \fa{لوازم} \fa{خانگی} \fa{میفروشد}.
\end{quote}

\begin{quote}
\emph{in forushgâh lavâzem-e khânegi miforushad.}
\end{quote}

\begin{quote}
This shop sells household appliances.
\end{quote}

\emph{lavâzem} 'things needed' (an Arabic plural) + \emph{khânegi} 'of the home'.

\begin{cousinweb}
\emph{lavâzem-e khânegi} is built in Persian on \emph{lavâzem}, from Arabic \emph{lawāzim}, 'necessities' (EtymDB), and \emph{khânegi}, 'of the home', from native \emph{khâne}, 'house'.
\end{cousinweb}

\begin{grammarlens}[title={What you've built}]
One more word for telling, explaining and describing at B1.
\end{grammarlens}

\subsection*{Your turn}
\begin{itemize}
\raggedright
  \item \emph{Say it:} \emph{lavâzem-e khânegi}
  \item \emph{Say it:} \emph{lavâzem-e khânegi}, once more
  \item \emph{Say it:} a sentence of your own with \emph{lavâzem-e khânegi}
  \item \emph{Recall:} say the Persian for a tenant, then the Persian for a microwave, then say \emph{lavâzem-e khânegi} again
\end{itemize}

\begin{tcolorbox}[breakable,colback=teal!4,colframe=teal!35!black,title={Before you move on}]
What does \fa{لوازم} \fa{خانگی} mean? (\textquotedblleft{}Household appliances\textquotedblright{}.) Say it once more.
\end{tcolorbox}
\section[\texorpdfstring{\fa{گرم} \fa{کردن} (garm kardan)}{گرم کردن (garm kardan)}]{\fa{گرم} \fa{کردن} (garm kardan) — to heat, to warm up}

\label{lesson:FA-C538-garm-kardan}



\begin{quote}
Before the new one: say the Persian for a microwave, then the Persian for household appliances.
\end{quote}

\subsection*{You'll want to know: \fa{گرم} \fa{کردن}}
\textbf{\fa{گرم} \fa{کردن}} (\emph{garm kardan}) — \textquotedblleft{}to heat, to warm up\textquotedblright{}.

\begin{quote}
\fa{او} \fa{شیر} \fa{بچه} \fa{را} \fa{گرم} \fa{کرد}.
\end{quote}

\begin{quote}
\emph{u shir-e bachche râ garm kard.}
\end{quote}

\begin{quote}
She warmed the baby's milk.
\end{quote}

\emph{garm} 'warm' + \emph{kardan}. Players also \emph{garm mikonand} 'warm up' before a match.

\begin{cousinweb}
\emph{garm kardan} is built on two inherited words: \emph{garm}, 'warm', which EtymDB traces to a Proto-Indo-European root meaning 'to heat, to warm' (EtymDB), and \emph{kardan}, 'to do'.
\end{cousinweb}

\begin{grammarlens}[title={What you've built}]
One more word for telling, explaining and describing at B1.
\end{grammarlens}

\subsection*{Your turn}
\begin{itemize}
\raggedright
  \item \emph{Say it:} \emph{garm kardan}
  \item \emph{Say it:} \emph{garm kardan}, once more
  \item \emph{Say it:} a sentence of your own with \emph{garm kardan}
  \item \emph{Recall:} say the Persian for a microwave, then the Persian for household appliances, then say \emph{garm kardan} again
\end{itemize}

\begin{tcolorbox}[breakable,colback=teal!4,colframe=teal!35!black,title={Before you move on}]
What does \fa{گرم} \fa{کردن} mean? (\textquotedblleft{}To heat, to warm up\textquotedblright{}.) Say it once more.
\end{tcolorbox}
\section[\texorpdfstring{\fa{سرد} \fa{شدن} (sard shodan)}{سرد شدن (sard shodan)}]{\fa{سرد} \fa{شدن} (sard shodan) — to get cold, to cool down}

\label{lesson:FA-C538-sard-shodan}



\begin{quote}
Before the new one: say the Persian for household appliances, then the Persian for to heat.
\end{quote}

\subsection*{You'll want to know: \fa{سرد} \fa{شدن}}
\textbf{\fa{سرد} \fa{شدن}} (\emph{sard shodan}) — \textquotedblleft{}to get cold, to cool down\textquotedblright{}.

\begin{quote}
\fa{چای} \fa{من} \fa{سرد} \fa{شد}.
\end{quote}

\begin{quote}
\emph{chây-e man sard shod.}
\end{quote}

\begin{quote}
My tea has gone cold.
\end{quote}

\emph{sard} 'cold' + \emph{shodan}. Also of feelings: \emph{ân do az ham sard shodand} 'those two grew cold to each other'.

\begin{cousinweb}
\emph{sard shodan} is built on two inherited words: \emph{sard}, 'cold', which EtymDB traces to Middle Persian (EtymDB), and \emph{shodan}, 'to become'.
\end{cousinweb}

\begin{grammarlens}[title={What you've built}]
One more word for telling, explaining and describing at B1.
\end{grammarlens}

\subsection*{Your turn}
\begin{itemize}
\raggedright
  \item \emph{Say it:} \emph{sard shodan}
  \item \emph{Say it:} \emph{sard shodan}, once more
  \item \emph{Say it:} a sentence of your own with \emph{sard shodan}
  \item \emph{Recall:} say the Persian for household appliances, then the Persian for to heat, then say \emph{sard shodan} again
\end{itemize}

\begin{tcolorbox}[breakable,colback=teal!4,colframe=teal!35!black,title={Before you move on}]
What does \fa{سرد} \fa{شدن} mean? (\textquotedblleft{}To get cold, to cool down\textquotedblright{}.) Say it once more.
\end{tcolorbox}
